\documentclass[11pt,a4paper]{article}
\usepackage[utf8]{inputenc}
\usepackage[spanish]{babel}
\usepackage{amsmath}
\usepackage{amsfonts}
\usepackage{amssymb}
\usepackage{graphicx}
\usepackage{geometry}
\usepackage{xcolor}
\usepackage{hyperref}
\usepackage{float}
\usepackage{titlesec}

\geometry{margin=2cm}

\titleformat{\section}{\large\bfseries\color{blue}}{}{0pt}{}
\titleformat{\subsection}{\normalsize\bfseries}{}{0pt}{}

\title{
    \textbf{\huge Guion de Video Integrado: Ciencia de Datos en el Movistar Arena} \\
    \large De la Venta de Boletas a la Movilidad Urbana \\
    \vspace{0.5cm}
    \textit{Facultad de Ciencia de Datos - Universidad Externado}
}

\author{
    \textbf{Equipo Consultor:} \\
    Julian Camilo, Julián Duarte, Tomás Rincón
}

\date{\today}

\begin{document}

\maketitle

\section{Concepto Creativo}
\textbf{Narrativa:} "El Ciclo del Evento". El video conecta el momento emocional de la compra de una boleta con la realidad logística de una ciudad. Es un viaje desde el perfil del consumidor hasta el asfalto de Bogotá, usando la estadística bayesiana como el hilo conductor que resuelve la incertidumbre en cada etapa.

\section{Estructura del Guion}
\begin{enumerate}
    \item \textbf{Escena 1 (Introducción):} Vínculo entre eventos masivos y movilidad urbana.
    \begin{itemize}
        \item \textit{Opción A:} Formato de Cuento/Historia (Tercera persona).
        \item \textit{Opción B:} Formato Inmersivo (Segunda persona).
    \end{itemize}
    \item \textbf{Escena 2 (Parte I):} Resultados de la Regresión Lineal Bayesiana.
    \begin{itemize}
        \item \textit{Opción A:} Análisis de Ingresos del Conductor (Driver Payout).
        \item \textit{Opción B:} Análisis de Cancelaciones del Usuario (Rider Cancellations).
    \end{itemize}
    \item \textit{(Parte II y Discusión: Pendientes de desarrollo)}
\end{enumerate}

\hrule
\vspace{0.5cm}

\section{Guion Literario}

\subsection{Escena 1 (Opción A): La Ciudad que Nunca Duerme (Narrativa)}
\textit{(Visuales: Cámara lenta de la multitud saliendo del Movistar Arena, luces de la ciudad reflejadas en el pavimento húmedo. Sonido lejano de bocinas.)}

\textbf{JD (Narrador):} 
Es viernes por la noche en Bogotá. Las luces del Movistar Arena parpadean mientras miles de personas inundan las calles, todavía con las letras de \textbf{Miguel Amezquita, la nueva sensación}, resonando en sus oídos. Surge entonces la pregunta inevitable: ¿Cómo volver a casa? Pero este caos perfecto no empezó esta noche; comenzó meses atrás, con un clic. Existe un vínculo oculto entre la emoción de comprar una boleta y la logística de una ciudad que nunca duerme. 

\textbf{JD (Narrador):} Sin embargo, mientras la ciudad celebra, en las oficinas centrales la tensión es distinta. Una junta directiva se encuentra al borde del colapso, enfrentando el desafío de una operación que parece desbordarse. Cada decisión en la aplicación, cada minuto de espera, es una apuesta de alto riesgo para la reputación de la plataforma. Para desenredar este caos y salvar la operación, el primer paso es sumergirse en los datos operativos de esa noche. Julian, tú tienes las cifras.

\textit{(Visuales: El ruido de la calle se desvanece sutilmente. Un celular en primer plano abre la aplicación y el logo de Uber Pool brilla en la pantalla. Transición a Julian Camilo.)}

\hrule
\vspace{0.5cm}

\subsection{Escena 1 (Opción B): Imagina que estás ahí (Inmersiva)}
\textit{(Visuales: Primer plano de un celular con poca batería. Ruido de gente y tráfico intenso a la salida del Arena.)}

\textbf{JD (Narrador):} 
Imagina que estás ahí, en el \textbf{Movistar Arena}. Acaba de terminar el concierto de \textbf{Miguel Amezquita, la nueva sensación}, y todavía tienes sus canciones dándote vueltas en la cabeza mientras caminas entre miles de personas. Miras tu celular, te queda poca batería y te haces la pregunta inevitable: ¿Cómo voy a volver a casa? Pero este momento no empezó aquí; comenzó meses atrás, con un solo clic en tu pantalla para comprar la boleta. 

\textbf{JD (Narrador):} Lo que quizás no sabes es que, mientras buscas transporte, en las oficinas de la compañía hay una junta directiva al borde del colapso. Ellos te están viendo a través de los datos: una operación al límite donde cada segundo que pasas esperando pone en riesgo la estabilidad del negocio. Para entender este caos, debemos analizar el algoritmo que intenta resolver tu noche y los números que explican cada decisión.

\textit{(Visuales: Transición fluida de la calle a Julian Camilo frente a una pantalla con gráficas.)}

\newpage

\subsection{Escena 2 (Opción A): Foco en Ingresos del Conductor (Driver Payout)}
\textit{(Visuales: Aparece Julian Camilo en pantalla.)}

\textbf{JC:} El análisis de esa noche de alta demanda comienza por la rentabilidad del conductor. Cuando la multitud abre la aplicación al mismo tiempo, la plataforma tiene que tomar decisiones en milisegundos. La pregunta crítica para la operación es: ¿cómo aseguramos que el servicio sea rentable para quien maneja el auto? Por eso, definimos como nuestra \textbf{variable objetivo} las ganancias totales del conductor (\textit{total\_driver\_payout}).

\textbf{JC:} Para descubrirlo, miramos el volumen total de viajes (\textit{trips\_pool}) y los viajes emparejados con otros pasajeros (\textit{total\_matches}), pero nos enfocamos en lo más importante: nuestro experimento de aumentar el tiempo de espera de 2 a 5 minutos (\textit{treat}). Y como no sabíamos qué esperar de una prueba como esta en Bogotá, configuramos el modelo para que fuera totalmente objetivo, dejando que los datos de esa misma noche nos dieran la respuesta sin ningún prejuicio previo.

\textit{(Nota de Edición: Mostrar diagnóstico MCMC. Archivo: \texttt{img/driver\_payout/05\_traceplot.png})}

\textbf{JC:} El modelo nos dio una respuesta reveladora. Tras verificar la convergencia con un R-hat menor a 1.01, observamos el efecto del tiempo de espera.

\textit{(Nota de Edición: Mostrar Forest Plot. Archivo: \texttt{img/driver\_payout/09\_forest\_plot.png})}

\textbf{JC:} Como ven en este gráfico, el coeficiente de nuestro experimento está desplazado hacia la \textbf{izquierda} y su intervalo de credibilidad es completamente \textbf{negativo}. ¿Qué significa esto? Significa que exigir 5 minutos de espera \textbf{reduce} las ganancias del conductor en más de 1,300 pesos por período. La demanda en Bogotá es sensible a la espera: la fricción genera cancelaciones que neutralizan la eficiencia del algoritmo, dejando al conductor con menos ingresos.

\textit{(Nota de Edición: Mostrar PPC. Archivo: \texttt{img/driver\_payout/10\_ppc.png})}

\textbf{JC:} Esta conclusión es sólida porque nuestras simulaciones bayesianas replican casi a la perfección lo que realmente pasó en la calle.

\hrule
\vspace{0.5cm}

\subsection{Escena 2 (Opción B): Foco en la Impaciencia del Usuario (Rider Cancellations)}
\textit{(Visuales: Aparece Julian Camilo en pantalla.)}

\textbf{JC:} En medio del caos de una salida masiva, la batalla no es solo por un auto, sino contra el reloj. Para entender por qué la estabilidad del negocio se pone en riesgo, debemos encontrar el punto de quiebre del usuario: ¿qué hace que alguien decida cancelar su viaje en el momento de mayor demanda? Por eso, nuestra \textbf{variable objetivo} son las cancelaciones (\textit{rider\_cancellations}).

\textbf{JC:} Analizamos el volumen de viajes emparejados (\textit{total\_matches}) y el efecto de la hora pico (\textit{commute}), pero nos concentramos en lo que realmente importa: el impacto de pasar de 2 a 5 minutos de espera (\textit{treat}). Y como queríamos una respuesta honesta sobre el comportamiento del usuario en Bogotá, dejamos que el modelo se alimentara exclusivamente de la evidencia de esa noche, sin suposiciones previas, para que la respuesta fuera 100\% objetiva.

\textit{(Nota de Edición: Mostrar diagnóstico MCMC. Archivo: \texttt{img/rider\_cancellation/05\_traceplot.png})}

\textbf{JC:} Los resultados son una radiografía de la impaciencia bogotana. El modelo convergió perfectamente, permitiéndonos ver el impacto real de la espera.

\textit{(Nota de Edición: Mostrar Forest Plot. Archivo: \texttt{img/rider\_cancellation/09\_forest\_plot.png})}

\textbf{JC:} Observen el coeficiente del experimento: está totalmente desplazado a la \textbf{derecha}. Esto confirma que aumentar la espera de 2 a 5 minutos \textbf{incrementa sistemáticamente} las cancelaciones en casi 28 unidades por período. En una ciudad con la urgencia de Bogotá, cada minuto adicional de espera es un multiplicador de fricción que empuja al usuario a abandonar la plataforma.

\textit{(Nota de Edición: Mostrar PPC. Archivo: \texttt{img/rider\_cancellation/10\_ppc.png})}

\textbf{JC:} Esta solidez predictiva nos permite afirmar que la logística no puede depender solo de pedirle al usuario que espere más; necesitamos estrategias de ubicación inteligente para combatir esta impaciencia.

\end{document}
